% ===================================================================
%  TAULUKKO N:o 2 — Ihmisruumis. — Välitunti. — Leikit.
%  Traduction 2026 d'apres texto/io, controlee sur texto/fr et
%  texto/sv. Consigne : voir texto/fi/10-taulukko-01.tex.
% ===================================================================

%%K t02-apar-1 apar
\VUsaut{4.0mm}
\VUtitre{15.0pt}{60}{TAULUKKO N:o 2}
\VUsaut{2.0mm}
\VUtitre{11.4pt}{0}{Ihmisruumis. --- Välitunti.}
\VUtitre{11.4pt}{0}{Leikit.}

%%K t02-tit-0 sub
\VUtitre{13.2pt}{0}{Ihmisruumis.}
\VUsaut{2.4mm}
\VUcentre{10.2pt}{0}{\textit{(Yksinkertainen luonnontiedon oppitunti.)}}
\VUsaut{3.0mm}

%%K t02-tit-1 sub pos=t02-01-1
\VUpk{10.2pt}{40}{Pää.}
\VUsaut{3.0mm}

%%K t02-01-1 p
1. --- Ihmisruumiin tärkeimmät osat ovat: \VUgras{pää}, \VUgras{vartalo} ja
\VUgras{raajat}.

%%K t02-02-1 p
2. --- Pää käsittää kaksi osaa: \VUgras{kallon}, joka sisältää \VUgras{aivot}
ja jota \VUgras{hiukset} peittävät, ja \VUgras{kasvot}.

%%K t02-03-1 p
3. --- Piirroksessamme emme näe kalloa emmekä aivoja; mutta näemme hyvin
selvästi kasvojen tärkeät osat.

%%K t02-04-1 p
4. --- Tässä on ensin \VUgras{otsa}, joka todistaa voimakkaasta ymmärryksestä,
lakkaamatta toimivasta ajattelusta. \VUgras{Ohimot} ovat hyvin vapaat.
\VUgras{Silmä} on eloisa ja selkoselällään. Tuuhea \VUgras{kulmakarva} ja pitkät
\VUgras{silmäripset} suojaavat sitä.

%%K t02-05-1 p
5. --- Tässä ovat lisäksi \VUgras{nenä} ja \VUgras{sieraimet}. Nenä palvelee
meitä haistamiseen ja hengittämiseen. Nenän kummallakin puolella ovat
\VUgras{posket}, kauempana \VUgras{korvat}. Kasvojen alaosan muodostavat
\VUgras{suu} ja \VUgras{leuka}.

%%K t02-06-1 p
6. --- Emme näe niitä kovin hyvin, sillä \VUgras{viikset} ja
\VUgras{pulisongit} peittävät ne meiltä. Mutta tiedämme, että suussa ovat kaksi
\VUgras{huulta}, \VUgras{yläleuka} ja \VUgras{alaleuka} \VUgras{hampaineen},
sekä \VUgras{kieli}. Suulla puhumme ja syömme. Kielellä maistamme ravinnon.

%%K t02-07-1 p
7. --- Silmä, korva, nenä ja suu ovat, samoin kuin sormet, viiden aistin
elimiä: \VUgras{näkö}, \VUgras{kuulo}, \VUgras{hajuaisti}, \VUgras{maku},
\VUgras{tunto}.

%%K t02-08-1 p
8. --- Pää on siis yksi ihmisruumiin mielenkiintoisimmista osista.

%%K t02-tit-2 sub
\VUsaut{4.4mm}
\VUpk{10.2pt}{40}{Vartalo.}
\VUsaut{3.0mm}

%%K t02-09-1 p
9. --- \VUgras{Kaula} yhdistää pään vartaloon. Se käsittää \VUgras{kurkun} ja
\VUgras{niskan}.

%%K t02-10-1 p
10. --- Vartalo käsittää myös monta olennaista elintä. \VUgras{Rinnassa} ovat
\VUgras{keuhkot}, joilla hengitämme, ja \VUgras{sydän}, joka on elämän
keskipiste. Kun sydän lakkaa lyömästä, elävä olento kuolee.

%%K t02-11-1 p
11. --- \VUgras{Kylkiluut} pitävät rintaa koossa.

%%K t02-12-1 p
12. --- Vartalon muut osat ovat \VUgras{kylki}, \VUgras{selkä},
\VUgras{olkapäät} ja \VUgras{vatsa}. Käymme makuulle kyljelle tai selälle
nukkuaksemme, kannamme taakat olkapäällämme.

%%K t02-13-1 p
13. --- Vatsassa ovat \VUgras{maha} ja \VUgras{suolet}. Ne palvelevat meitä
ravinnon sulattamiseen.

%%K t02-tit-3 sub
\VUsaut{4.4mm}
\VUpk{10.2pt}{40}{Raajat.}
\VUsaut{3.0mm}

%%K t02-14-1 p
14. --- Meillä on neljä \VUgras{raajaa}: kaksi \VUgras{käsivartta} ja kaksi
\VUgras{jalkaa}. Käsivarsilla otamme, pidämme, nostamme ja kokoamme esineet;
jaloilla seisomme, kävelemme ja juoksemme.

%%K t02-15-1 p
15. --- \VUgras{Olkapää} yhdistää käsivarren ruumiiseen. Hyvin voimakas
\VUgras{lihas}, hauis, liikuttaa käsivartta, jonka \VUgras{kyynärpää} yhdistää
\VUgras{kyynärvarteen}. \VUgras{Ranne} muodostaa \VUgras{käden} nivelen, joka
päättyy \VUgras{sormiin} ja \VUgras{kynsiin}. Kädessä erotamme
\VUgras{laskimon} (ja valtimon), jossa \VUgras{veri} virtaa.

%%K t02-16-1 p
16. --- Jalan osat ovat: \VUgras{reisi}, \VUgras{polvi} ja
\VUgras{polvitaive}, \VUgras{pohje}, jonka \VUgras{lihaa} \VUgras{iho} peittää,
\VUgras{nilkka} ja \VUgras{jalkaterä}. Jalan \VUgras{luut} ovat hyvin karkeat ja
hyvin vahvat. Jalkaterän päässä ovat \VUgras{varpaat}. Muut osat ovat
\VUgras{jalkapohja} ja \VUgras{kantapää}.

%%K t02-17-1 p
17. --- Sellainen on ihmisruumiin miltei täydellinen kuvaus.

%%K t02-17-2 p
\textit{Huom.} --- \textit{Ilmauksen: « minulla on kipeä... (pää j. n. e.) »
avulla opettaja voi käyttää koko tämän sanavaraston uudelleen ruumiin hyvän tai
huonon} \VUgras{tilan} \textit{kannalta.}

%%K t02-tit-4 sub
\VUsaut{5.0mm}
\VUpk{10.2pt}{40}{Voimistelu.}
\VUsaut{2.4mm}
\VUcentre{10.2pt}{0}{\textit{(Erään oppilaan kertomus.)}}
\VUsaut{3.0mm}

%%K t02-18-1 p
18. --- Ollaksemme notkeita, ketteriä ja terveitä meidän täytyy harjoittaa
jalkojamme ja käsivarsiamme. Siksi leikimme ja harrastamme \VUgras{voimistelua}.

%%K t02-19-1 p
19. --- Viikolla nämä kaksi harjoitusta tapahtuvat oppilaitoksen pihalla (katso
taulukko n:o 1). Mutta torstaisin ja sunnuntaisin menemme suurelle
\VUgras{nurmikentälle}, jossa meillä on tilaa juosta ja huvitella.

%%K t02-20-1 p
20. --- Opettajamme ja vanhempamme seuraavat meitä joskus sinne; mutta
\VUgras{voimistelunopettaja} \textsuperscript{(12)} tietenkin johtaa
harjoituksia.

%%K t02-21-1 p
21. --- Nurmikentällä, sisäänkäynnin vasemmalla puolella, ovat
\VUgras{voimisteluvälineet} \textsuperscript{(1)}, kiipeämme \VUgras{tikkaita}
\textsuperscript{(2)} ylös, riipumme pää alaspäin \VUgras{renkaissa}
\textsuperscript{(3)} ja \VUgras{trapetsissa} \textsuperscript{(4)}, vedämme
itsemme käsin \VUgras{köysitikkaiden} \textsuperscript{(5)} tai
\VUgras{solmuköyden} \textsuperscript{(6)} huipulle.

%%K t02-22-1 p
22. --- Sillä aikaa toverimme hyppäävät \VUgras{puuhevosen}
\textsuperscript{(7)} tai \VUgras{nojapuiden} \textsuperscript{(11)} yli. Toiset
\VUgras{nyrkkeilevät} \textsuperscript{(8)} tai \VUgras{miekkailevat}
\textsuperscript{(9)} naamio kasvoilla ja \VUgras{floretti} kädessä. Lopuksi
toiset \VUgras{nostavat} \VUgras{levytankoja} \textsuperscript{(10)}.

%%K t02-tit-5 sub
\VUsaut{4.6mm}
\VUpk{10.2pt}{40}{Leikit.}
\VUsaut{3.0mm}

%%K t02-23-1 p
23. --- Voimistelijoiden ympärillä pojat ja tytöt leikkivät erilaisia
\VUgras{leikkejä}. Tytöt huvittelevat \VUgras{vanteellaan}
\textsuperscript{(14)}; he hyppäävät \VUgras{narua} \textsuperscript{(13)} tai
kutsuvat veljiään ja serkkujaan \VUgras{kroketti}-otteluun
\textsuperscript{(15)}. Heitä ilahduttaa työntää pois vastustajan
\VUgras{pallo} \VUgras{puunuijallaan} juuri sillä hetkellä, kun tämä luulee
menevänsä helposti portin läpi!

%%K t02-24-1 p
24. --- Pojat suosivat leikkejä, jotka vaativat enemmän lihaksia. Mielellään he
leikkivät \VUgras{keiloilla} \textsuperscript{(16)}, jotka he kaatavat taitavasti
\VUgras{pallolla} \textsuperscript{(17)}. He eivät myöskään halveksi leikkiä
\VUgras{hyrrällä} \textsuperscript{(18)} ja \VUgras{bilboketilla}
\textsuperscript{(19)} tai jopa \VUgras{sammakkoleikkiä} \textsuperscript{(20)}.
Mutta heidän lempileikkejään ovat tietenkin \VUgras{kuulat}
\textsuperscript{(21)} ja \VUgras{seinäpallo} \textsuperscript{(22)}, sillä
\VUgras{kasvihuoneen} \textsuperscript{(26)} seinä on hyvin sopiva
kimmottamaan kevyttä palloa takaisin.

%%K t02-25-1 p
25. --- Pikku Juhani, joka kävelee \VUgras{puujaloillaan}
\textsuperscript{(24)}, on hyvin taitava ja osaa hyvin säilyttää tasapainonsa.
Hänen lähellään muutamat toverit leikkivät \VUgras{piilosta}
\textsuperscript{(25)} tuossa kasvihuoneessa ja \VUgras{pensasaidan} takana.
Toiset suosivat \VUgras{lyönninarvausta} \textsuperscript{(28)} tai
\VUgras{sulkapalloa} \textsuperscript{(29)}, joka lentelee \VUgras{mailasta}
toiseen.

%%K t02-noto-1 noto
\VUnotes{143.2mm}{%
(*) Lasten- tai poikienleikeillä on usein puhtaasti kansallisia nimiä. Olemme
toimineet nimetäksemme ne idoksi niin hyvin kuin mahdollista, joko antamalla
itsemme innoittua itse teosta, joka niille on ominaista, tai valitsemalla
kansallisen nimen siinä tai tässä maassa, kun se ei ole liian väärä. Esim.:
\VUgras{tynnyrileikki} (saksaksi ja ranskaksi) on \VUgras{sammakkoleikki}
\textsuperscript{(20)} (englanniksi), jossa leikkijät koettavat heittää
pyöreitä kiekkojaan sen metallisammakon suun läpi, joka suu ammollaan seisoo
lävistetyn laatikon (tynnyrin) päällä. \VUgras{Olkahyppy}
\textsuperscript{(30)} eroaa \VUgras{selkähypystä} vain tässä: pään yli
hypätäkseen leikkijät nojaavat kädet toverinsa olkapäihin, kun taas
« \VUgras{selkähypyssä} » he nojaavat kädet vain hänen selkäänsä. --- Tai myös
jotkut osanottajista ovat kumarassa sillä aikaa kun toiset hyppäävät heidän
selkänsä yli nojaten kädet olkapäähän; ensimmäisten täytyy kestää isku
kyyristymättä, toisten täytyy hypätä tasapainoa menettämättä (Katso itse
taulukkoa).%
}

%%K t02-26-1 p
26. --- Keskellä nurmikenttää on kaksi puuta. Toisen näiden kahden juurella
seitsemän tai kahdeksan poikaa on järjestänyt \VUgras{olkahyppy}-ottelun
\textsuperscript{(30)} (*). Kaikissa maissa ei tunneta tuota leikkiä; mutta
kaikki tietävät, mitä \VUgras{selkähyppy} on, jota pojat eivät tällä kertaa
leiki.

%%K t02-tit-6 sub
\VUsaut{5.0mm}
\VUpk{10.2pt}{40}{Leikit ulkoilmassa.}
\VUsaut{3.4mm}

%%K t02-27-1 p
27. \VUgras{Sokkosilla} \textsuperscript{(31)} on myös hyvin hauskaa; tytöt ja
pojat panevat vuorotellen \VUgras{siteen} silmilleen ja tarttuvat
kömpelöihin, jotka antavat pyydystää itsensä.

%%K t02-28-1 p
28. --- Keskellä nurmikenttää, tässä on hyvin ranskalainen mutta hieman
väkivaltainen leikki: se on \VUgras{karhuemo} \textsuperscript{(32)}, joka on
paljon meluisampi kuin niin rauhallinen leikki \VUgras{leijalla}
\textsuperscript{(33)} tai jopa kuin englantilainen leikki \VUgras{tennis}
\textsuperscript{(34)}.

%%K t02-29-1 p
29. --- Viimeksi mainittu urheilu on isojen oppilaiden ilo. Opettajamme itse
harrastavat sitä mielellään. Se vaatii suurta taitoa ja ketteryyttä, ja se
miellyttää minua kovasti. Jätän tytöille leikit \VUgras{keinulla}
\textsuperscript{(35)}.

%%K t02-30-1 p
30. --- En pidä eräästä toisesta englantilaisesta leikistä, josta ehkä tulisi
vaarallinen.

%%K t02-30-1b p
Tarkoitan \VUgras{jalkapalloa} \textsuperscript{(36)}, joka ilahduttaa
vanhempaa veljeäni. Suosin vanhaa \VUgras{puomileikkiämme}
\textsuperscript{(38)}, joka on yhtä terveellinen ja todella vähemmän raaka.

%%K t02-tit-7 sub
\VUsaut{4.6mm}
\VUpk{10.2pt}{40}{Pyöräily ja pallopelit.}
\VUsaut{3.0mm}

%%K t02-31-1 p
31. --- Myös mielelläni ajan nurmikentän ympäri \VUgras{polkupyörällä}
\textsuperscript{(37)}.

%%K t02-32-1 p
32. --- Ensi vuonna opettelen \VUgras{ratsastuksen} \textsuperscript{(39)}.
Kuinka kaunis onkaan hevonen, joka sulavasti astelee tai ravaa, ja joka
laukassa hyppää esteiden yli!

%%K t02-33-1 p
33. --- Kesällä opimme \VUgras{uimataidon} \textsuperscript{(40)}. Sukellamme
ihastuksissamme ja kylvemme joen kirkkaassa ja raikkaassa vedessä. Joskus
lopuksi päästämme ilmaan paperisen \VUgras{ilmapallon} \textsuperscript{(41)},
joka majesteettisesti nousee ilmaan heti kun se on täytetty lämpimällä ilmalla.

%%K t02-34-1 p
34. --- On myös monta muuta leikkiä, mutta en koskaan lakkaisi luettelemasta
niitä, ja tästä yhteenvedosta ymmärretään, ettei torstaisin ole kovin ikävää
kauniilla nurmikentällämme.

%%K t02-34-2 p
Huom. --- \textit{Opettaja voi helposti täydentää tuon luvun; kuvaamalla
yksityiskohtaisemmin kunkin tärkeimmistä leikeistä.}
\VUsaut{6.0mm}
\VUfilet{22mm}
